\begin{tikzpicture}[font=\small,>=stealth,line join=round]
\draw[gray!40,rounded corners=3pt] (0,0) rectangle (14.2,9.4);
\node at (7.1,8.8) {分子分母成比例：点在移动，连线斜率却保持不变};
\begin{scope}[shift={(3.3,1.8)},x=.75cm,y=.75cm]
\draw[->,gray] (0,0)--(3.4,0) node[right]{$X$};\draw[->,gray] (0,0)--(0,7.5) node[above]{$Y$};
\draw[blue,very thick,->] (1,2)--(3.3,6.6);
\fill[blue] (1,2) circle(2.2pt) node[above left=5pt]{$(1,2)$};
\node[below left=4pt] at (0,0){$O$};
\end{scope}
\node[align=left,anchor=west] at (8,5.1){$X=x^2+1\ge1$\\[9pt]$Y=2x^2+2=2X$\\[12pt]动点沿射线运动\\[8pt]原点连线斜率始终是 $2$\\[12pt]轨迹不是抛物线\\[8pt]函数值也不随点移动而变};
\node[text width=13cm,align=center] at (7.1,.6) {横纵坐标等比例；当 $d\ne0$ 时横坐标会变化，轨迹不会固定成一个点。};
\end{tikzpicture}
